
\textup{D)} Suppose first that \(S\) is the spectrum of a field.
If \(F'(0)=0\), then \(F=1\). Formula \((5'')\) then gives
\[
  H(x+y)=H(x)H(y),
\]
which shows that \(H\) extends to a group morphism
\[
  \mathbb G_{a,S}\longrightarrow T
\]
\textup{(\SGARefLink{sga3:xviii:prop:2.3}{Exposé XVIII, 2.3})}, necessarily the constant morphism
with value \(e\). On the other hand, \SGARefLink{sga3:xx:2.3}{Lemma 2.3} also gives
\[
  E(x)=1+bx
\]
for some \(b\in k\). Formula \((6'')\) therefore gives, for
\((x,y)\in N(S')\),
\[
  (x+y)E(x+y)^{-1}=xE(x)^{-1}+yE(y)^{-1}.
\]
Thus,\footnote{The editors added the sentence that follows. One
may also see by direct calculation that the preceding equality
implies
\(0=xyb\bigl(2+(x+y)b\bigr)\), hence
\(0=b\bigl(2+(x+y)b\bigr)\), and finally \(b=0\).}
again by \SGARefLink{sga3:xviii:prop:2.3}{Exposé XVIII, 2.3}, the map
\[
  x\longmapsto xE(x)^{-1}
\]
extends to a morphism of \(k\)-groups
\(\mathbb G_{a,k}\to\mathbb G_{a,k}\). Consequently
\[
  \frac{x}{1+bx}=cx
\]
for some \(c\in k\), whence \(b=0\) and \(c=1\).

Thus \(F,H,E\) are constant with values \((1,e,1)\) in a
neighborhood of the zero section, and hence everywhere. Formula
\((+)\) then shows that \(U_\alpha\) and \(U_{-\alpha}\) commute,
contrary to hypothesis \textup{\SGARefLink{sga3:xx:2.3:item:ii}{(ii)}}.

For arbitrary \(S\), we have therefore proved that \(F'(0)\) is
nonzero on every fiber, and hence invertible. The same is plainly
true of \(E'(0)\). \SGARefLink{sga3:xx:2.3}{Lemma 2.3} now shows that there are
\(a,b\in\mathbb G_m(S)\) such that
\SGARefTarget{sga3:xx:2.2:diamond-one}
\begin{equation}\tag{\(\diamondsuit_1\)}
  F(x)=1+ax,\qquad E(x)=1+bx
  \qquad\text{for }x\in W(S').
\end{equation}

\textup{E)} The rest is now easy. Substitution of the preceding
results in \((3'')\) gives
\[
  y\,\alpha\bigl(H(x)\bigr)(1+ay)
   =y\bigl(1+ax+ay(1+bx)\bigr)(1+ax).
\]
This formula holds for every section \((x,y)\) of \(N\). Since
\(\mathbb G_{a,S}\times_S\mathbb G_{m,S}\) is schematically dense
in \(\mathbb G_{a,S}^{\,2}\), it follows that
\[
  (1+ay)\alpha\bigl(H(x)\bigr)
   =\bigl(1+ax+ay(1+bx)\bigr)(1+ax).
\]
Setting \(y=0\) gives
\[
  \alpha\bigl(H(x)\bigr)=(1+ax)^2.
\]
Substitution in the preceding equality gives\footnote{In the three
preceding equalities, the editors corrected the original by
replacing the factor \(1+bx\) on the right by \(1+ax\).
Substitution of the third equality in the second, and use of the
invertibility of \(1+ax\), gives precisely
\(a^2xy=abxy\).}
\[
  a^2xy=abxy.
\]
Since \(\mathbb G_{m,S}\) is schematically dense in
\(\mathbb G_{a,S}\), we obtain \(a^2=ab\). As \(a\) is invertible,
\SGARefTarget{sga3:xx:2.2:diamond-two}
\begin{equation}\tag{\(\diamondsuit_2\)}
  a=b,\qquad \alpha\bigl(H(x)\bigr)=(1+ax)^2.
\end{equation}

Because \(a\) is invertible, the map \(x\mapsto1+ax\) is an
automorphism of \(\mathbb G_{a,S}\). There are therefore an open
subset \(W'\) of \(\mathbb G_{a,S}\), containing the section \(1\),
and a morphism
\[
  P\colon W'\longrightarrow T
\]
such that
\[
  P(1+ax)=H(x).
\]
\footnote{This is the change of variables
\(x'=1+ax\), or \(x=(x'-1)/a\).}
Substitution in \((2')\) gives immediately, for
\((x,y)\in N(S')\),
\[
  P(1+ax+ay)
   =P\left(\frac{1+ax+ay}{1+ax}\right)P(1+ax).
\]
Thus there is an open neighborhood of \(1\) in
\(\mathbb G_{m,S}\) on which
\[
  P(x)P(y)=P(xy).
\]
By \SGARefLink{sga3:xviii:prop:2.3}{Exposé XVIII, 2.3}, there is a group morphism
\SGARefTarget{sga3:xx:2.2:diamond-three}
\begin{equation}\tag{\(\diamondsuit_3\)}
  \alpha^\ast\colon\mathbb G_{m,S}\longrightarrow T
\end{equation}
extending \(P\). Since
\(\alpha(H(x))=(1+ax)^2\) near the zero section, we have
\(\alpha(\alpha^\ast(z))=z^2\) near the section \(1\), and hence
\SGARefTarget{sga3:xx:2.2:diamond-four}
\begin{equation}\tag{\(\diamondsuit_4\)}
  \alpha\circ\alpha^\ast=2.
\end{equation}

\textup{F)}\footnote{The editors slightly modified what follows,
since an open subset \(V\) had already been introduced in part A.}
Combining \((+)\) and
\((\diamondsuit_1)\)--\((\diamondsuit_4)\), we see that there are
\[
  a\in\mathbb G_m(S),\qquad
  \alpha^\ast\in
    \operatorname{Hom}_{S\text{-gr.}}(\mathbb G_{m,S},T)
\]
such that \(\alpha\circ\alpha^\ast=2\), and such that, whenever
\((u,v)\in V(S')\) and \(uv\in W(S')\), the element \(1+auv\)
is invertible and
\[
  p_\alpha(u)p_{-\alpha}(v)
   =p_{-\alpha}\left(\frac{v}{1+auv}\right)
      \alpha^\ast(1+auv)
      p_\alpha\left(\frac{u}{1+auv}\right).
\]

Let \(V'\) be the open subset of \(\mathbb G_{a,S}^{\,2}\) defined
by the invertibility of \(1+auv\); thus
\[
  V'=\bigl(\mathbb G_{a,S}^{\,2}\bigr)_f,
  \qquad f(u,v)=1+auv.
\]
The two sides of the preceding formula define morphisms
\(V'\to G\) which agree in a neighborhood of the zero section,
and hence agree on \(V'\). Thus the formula holds for every
section \((u,v)\) of \(V'\), and in particular \(V'\subset U\),
where \(U\) is the open subset introduced at the beginning of
part A.

We prove that \(U=V'\). Returning to the notation of part A, there
is a morphism
\[
  A\colon U\longrightarrow\mathbb G_{a,S}
\]
whose restriction to \(V'\) is
\[
  A(u,v)=v(1+auv)^{-1}.
\]
The equality \(U=V'\) is a set-theoretic question and may therefore
be checked when \(S=\operatorname{Spec}(k)\) for a field \(k\).
It then reduces to the evident assertion that the domain of
definition of the rational map
\[
  \mathbb G_{a,k}^{\,2}\dashrightarrow\mathbb G_{a,k},
  \qquad
  (X,Y)\longmapsto\frac{Y}{1+aXY},
\]
is the open subset defined by \(1+aXY\).

\textup{G)} We have proved the existence of \(a\) and
\(\alpha^\ast\), as well as the two asserted supplementary
properties. It remains to prove uniqueness. Let
\(a'\) and \(\alpha^{\ast\prime}\) also satisfy the required
conditions. For \(u,v\in\mathbb G_a(S')\), we immediately have
\[
\begin{gathered}
  1+auv\text{ invertible}
    \ \Longrightarrow\
  1+a'uv\text{ invertible},\\
  \frac{v}{1+auv}=\frac{v}{1+a'uv}.
\end{gathered}
\]
Hence, for every \(u\in\mathbb G_a(S')\),
\[
  1+au\text{ invertible}
    \ \Longrightarrow\
  1+au=1+a'u,
\]
which proves \(a=a'\). With the same notation,
\[
  1+au\text{ invertible}
    \ \Longrightarrow\
  \alpha^\ast(1+au)=\alpha^{\ast\prime}(1+au),
\]
and therefore \(\alpha^\ast=\alpha^{\ast\prime}\).
\qed
